% ==============================================================================
% PREAMBLE.TEX - Configuración, Paquetes y Entornos Estilizados
% ==============================================================================

\usepackage[spanish,es-tabla]{babel}
\usepackage[utf8]{inputenc}
\usepackage[margin=2.5cm, headheight=14pt]{geometry}
\raggedbottom % Evita errores 'Underfull \vbox' equilibrando márgenes inferiores

\usepackage{amsmath, amsfonts, amssymb, amsthm, mathtools}
\usepackage{enumitem}
\usepackage{booktabs}   % Tablas de calidad editorial
\usepackage{tabularx}   % Tablas ajustables
\usepackage{microtype}  % Mejora tipográfica de justificación
\usepackage{xcolor}
\usepackage[most]{tcolorbox}
\tcbuselibrary{theorems, skins, breakable}
\usepackage{hyperref}
\usepackage{bookmark}

\usepackage{tikz}
\usetikzlibrary{babel}
\usetikzlibrary{patterns}

% Registro completo de niveles en la tabla de contenidos
\setcounter{tocdepth}{2}
\setcounter{secnumdepth}{2}

\newcounter{problemacontador}[chapter]

\newcommand{\seccionproblemas}[1][Problemas resueltos]{%
  \section{#1}%
  \setcounter{problemacontador}{0}%
}


% ------------------------------------------------------------------------------
% Paleta de Colores Académicos
% ------------------------------------------------------------------------------
\definecolor{maincolor}{RGB}{27, 54, 93}      % Azul Marino Profundo (Títulos/Encabezados)
\definecolor{accentcolor}{RGB}{0, 128, 128}    % Azul Teal / Turquesa Oscuro
\definecolor{defcolor}{RGB}{41, 128, 185}     % Azul Definición
\definecolor{conceptcolor}{RGB}{31, 78, 121}   % Azul Concepto Clave
\definecolor{excolor}{RGB}{39, 174, 96}       % Verde Ejemplo Resuelto
\definecolor{notecolor}{RGB}{142, 68, 173}    % Morado Nota / Observación
\definecolor{bglight}{RGB}{248, 249, 250}     % Fondo neutro claro
\definecolor{boxbg}{RGB}{245, 247, 250}       % Fondo claro para cajas

% ------------------------------------------------------------------------------
% Configuración de Hyperref
% ------------------------------------------------------------------------------
\hypersetup{
    colorlinks=true,
    linkcolor=maincolor,
    citecolor=accentcolor,
    urlcolor=accentcolor,
    pdfstartview=FitH
}

% ------------------------------------------------------------------------------
% Configuración de Encabezados y Pies a Doble Cara (fancyhdr)
% ------------------------------------------------------------------------------
\usepackage{fancyhdr}
\pagestyle{fancy}
\fancyhf{} 

% Páginas Pares (Izquierda / Left): Capítulo a la izq, Número de página a la der.
\fancyhead[LE]{\small\bfseries\color{maincolor}\leftmark}
\fancyhead[RE]{\small\bfseries\color{maincolor}\thepage}

% Páginas Impares (Derecha / Right): Número de página a la izq, Sección a la der.
\fancyhead[LO]{\small\bfseries\color{maincolor}\thepage}
\fancyhead[RO]{\small\bfseries\color{maincolor}\rightmark}

\renewcommand{\chaptermark}[1]{\markboth{\thechapter.\ #1}{}}
\renewcommand{\sectionmark}[1]{\markright{\thesection.\ #1}}
\renewcommand{\headrulewidth}{0.4pt}

% Estilo plano para portadas e inicios de capítulo
\fancypagestyle{plain}{
  \fancyhf{}
  \fancyfoot[C]{\thepage}
  \renewcommand{\headrulewidth}{0pt}
}

% ------------------------------------------------------------------------------
% Entornos Estilizados (tcolorbox)
% ------------------------------------------------------------------------------

% 1. Entorno Definición
\newtcolorbox[auto counter, number within=chapter]{definicion}[1][]{%
  enhanced,
  breakable,
  colback=bglight,
  colframe=defcolor,
  fonttitle=\bfseries\color{white},
  title={Definición \thetcbcounter\if\relax\detokenize{#1}\relax\else\ (#1)\fi},
  attach boxed title to top left,
  boxed title style={colback=defcolor, sharp corners},
  arc=2pt,
  boxrule=0.8pt,
  top=8pt, bottom=8pt, left=10pt, right=10pt
}

% 2. Entorno Concepto / Principio Clave
\newtcolorbox[auto counter, number within=chapter]{concepto}[1][]{%
  enhanced,
  breakable,
  colback=maincolor!4!white,
  colframe=conceptcolor,
  fonttitle=\bfseries\color{white},
  title={Principio Fundamental \thetcbcounter\if\relax\detokenize{#1}\relax\else\ (#1)\fi},
  attach boxed title to top left,
  boxed title style={colback=conceptcolor, sharp corners},
  arc=2pt,
  boxrule=0.8pt,
  top=8pt, bottom=8pt, left=10pt, right=10pt
}

% 3. Entorno Ejemplo Resuelto
\newtcolorbox[auto counter, number within=chapter]{ejemplo}[1][]{%
  enhanced,
  breakable,
  colback=excolor!4!white,
  colframe=excolor,
  fonttitle=\bfseries\color{white},
  title={Ejemplo Resuelto \thetcbcounter\if\relax\detokenize{#1}\relax\else\ (#1)\fi},
  attach boxed title to top left,
  boxed title style={colback=excolor, sharp corners},
  arc=2pt,
  boxrule=0.8pt,
  top=8pt, bottom=8pt, left=10pt, right=10pt
}

% 4. Entorno Observación / Nota Técnica
\newtcolorbox{observacion}[1][Observación]{%
  enhanced,
  breakable,
  colback=notecolor!4!white,
  colframe=notecolor,
  fonttitle=\bfseries\color{white},
  title={#1},
  attach boxed title to top left,
  boxed title style={colback=notecolor, sharp corners},
  arc=2pt,
  boxrule=0.8pt,
  top=8pt, bottom=8pt, left=10pt, right=10pt
}

% ------------------------------------------------------------------------------
% Entorno para Problemas de la Lista del Profesor
% ------------------------------------------------------------------------------
% 8. PROBLEMA RESUELTO (Púrpura)
\newtcolorbox{problema}[1][]{
  enhanced,
  breakable,
  colback=boxbg,
  colframe=maincolor,
  fonttitle=\bfseries,
  coltitle=white,
  arc=2mm,
  boxrule=1pt,
  code={\refstepcounter{problemacontador}},
  title={Problema~\theproblemacontador\ifstrempty{#1}{}{: #1}},
  lower separated=true
}

% ------------------------------------------------------------------------------
% Macros y Comandos de Física de Partículas
% ------------------------------------------------------------------------------
\newcommand{\eV}{\text{ eV}}
\newcommand{\keV}{\text{ keV}}
\newcommand{\MeV}{\text{ MeV}}
\newcommand{\GeV}{\text{ GeV}}
\newcommand{\TeV}{\text{ TeV}}
\newcommand{\fm}{\text{ fm}}
\newcommand{\SU}[1]{\text{SU}(#1)}
\newcommand{\U}[1]{\text{U}(#1)}
\newcommand{\SO}[1]{\text{SO}(#1)}
\newcommand{\bra}[1]{\langle #1 |}
\newcommand{\ket}[1]{| #1 \rangle}
\newcommand{\bracket}[2]{\langle #1 | #2 \rangle}